\documentclass[9pt,twocolumn]{extarticle}
\usepackage[letterpaper,margin=0.43in,columnsep=0.18in]{geometry}
\usepackage{amsmath,amssymb,balance,booktabs,microtype,xcolor}
\usepackage[hidelinks]{hyperref}
\usepackage{enumitem}
\setlength{\parindent}{0pt}
\setlength{\parskip}{1.6pt}
\setlength{\abovecaptionskip}{2pt}
\setlength{\belowcaptionskip}{-2pt}
\setlength{\textfloatsep}{3pt}
\setlength{\floatsep}{3pt}
\setlist{nosep,leftmargin=9pt}
\newcommand{\state}{S}
\newcommand{\code}{C}

\title{\vspace{-1.35em}\textbf{Closed State Channels Turn Local Language-Model
Competence into Long-Horizon Computation}\vspace{-0.65em}}
\author{\normalsize Anonymous authors}
\date{}

\begin{document}
\maketitle
\vspace{-1.45em}
\begin{abstract}
\small
Reasoning language models can execute one state update yet fail to compose the
same update over a history.  We force Qwen models through a one-token discrete
state channel and delete used history after each call.  A two-bit channel hits
its exact 50\% information limit; direct transition-closure training makes a
redundant four-bit channel 93.25\% accurate at horizon 16 under five shifted
history distributions, versus 48.50\% for matched endpoint training.  Predicted
codes carry a stable three-bit semantic quotient despite surface-token drift.
\end{abstract}
\vspace{-0.55em}
\balance

\section*{\fontsize{9}{10}\selectfont Question and contract}
\vspace{-0.3em}
A history $H$ induces the sufficient current state
$\state=f(H)\in\{0,\ldots,7\}$; a held-out final rule $F$ gives
$Y=g(\state,F)$.  Native one-pass execution receives $(H,F)$.
Our interface instead emits one token $\code$, deletes $H$, and gives only
$(\code,F)$ to a second call.  Repeated execution also requires a learned
transition $\widehat T_u$ for each operation $u$:
\[
\underbrace{Y\perp H\mid(\code,F)}_{\text{sufficiency}},\qquad
\underbrace{q(\widehat T_u(c))=T_u(q(c))}_{\text{semantic closure}},\qquad
R=\log_2|\mathcal C|.
\]
Here $q$ decodes a token to state.  Exact canonicalization asks that equal
states emit the same token; the weaker and more useful requirement allows a
\emph{fiber} $q^{-1}(s)$ of equivalent tokens.  This is a finite, causal version
of a predictive sufficient statistic \cite{shalizi2001}, with a declared
channel rate rather than an unconstrained hidden vector
\cite{shannon1948,tishby2000}.

\section*{\fontsize{9}{10}\selectfont Method}
\vspace{-0.3em}
\textbf{Diagnosis.}  On 1,920 balanced modulo-eight programs, frozen
Qwen2.5-32B has perfect Read, one-step Update, and supplied-state use, but only
13.44\% h2 one-pass Compose.  A two-call self-handoff reaches 76.98\%; stepwise
decimal execution remains perfect through h4.

\textbf{Training.}  Rank-16 LoRA on Qwen2.5-7B uses group-disjoint program
contexts.  Each semantic program supplies two 256-token forwards and two
one-token targets: producer $(H\!\to\!\code)$ or
$(\code,u\!\to\!\code')$, and consumer $(\code,F\!\to\!Y)$.  Prompt tokens
carry no loss.  Endpoint controls use the same programs, forwards, target
count, padding, and optimizer, but train $(H\!\to\!\code)$ rather than closure.
Evaluation composes two-operation calls to h16, discarding each used block.

\textbf{Controls and interventions.}  We compare four codes for eight states
(2 bits), one opaque code per state (3 bits), and two tokens per state (4
bits).  Producer and consumer contexts are disjoint; test contexts are unseen.
We measure exact token accuracy, $I(\state;\code)$, closure after applying $q$,
same-state agreement before/after quotienting, and whether an independently
predicted donor code makes the recipient consumer follow the donor state.

\begin{center}
\centering
\scriptsize
\setlength{\tabcolsep}{2.5pt}
\begin{tabular}{@{}lrrl@{}}
\toprule
Interface / training & h8 & h16 & decisive check\\
\midrule
decimal recursive & 100.0 & 100.0 & also 100\% at h32\\
2-bit compressed & 50.0 & 50.0 & exact rate limit\\
3-bit endpoint & 35.0 & 24.5 & enough rate, no closure\\
3-bit closure & 55.0 & 55.0 & entry encoder damaged\\
4-bit endpoint & 54.0 & 52.5 & matched control\\
4-bit closure & \textbf{97.8} & \textbf{97.1} & $+43.9,+44.6$ points\\
\bottomrule
\end{tabular}
\par\vspace{0.15em}
\parbox{0.98\columnwidth}{\scriptsize\textbf{Table 1:}
Matched addition programs; accuracy (\%).  Context-paired
95\% intervals for the 4-bit closure gains are $[36.5,51.0]$ at h8 and
$[37.6,51.5]$ at h16.}
\end{center}

\section*{\fontsize{9}{10}\selectfont Findings and theory}
\vspace{-0.3em}
\textbf{Rate is causal.}  With uniform $\state$ and four deterministic code
classes, every token aliases two states, so an arbitrary one-to-one final table
permits at most $4/8=50\%$ exact recovery.  The trained two-bit channel achieves
50\% at every horizon while closing perfectly.

\textbf{Closure, not endpoint fit, produces reuse.}  Both 4-bit adapters fit
short training, but only transition supervision survives recursive use.  Across
structured, IID, shuffled, cancellation, and repeated histories, closure scores
97.75/95.25/95.00/93.25\% at h2/4/8/16; endpoint training scores
100/81.25/63.50/48.50\%.  At h16, every family remains 90--100\%.

\textbf{The state is a quotient, not a token.}  In the matched h16 run, only
2.92\% of 4-bit outputs have the wrong state, while 42.40\% change only the
extra bit.  Thus exact-token accuracy is 54.69\%, but quotient agreement is
95.10\%, $I(\state;\code)=2.83$ bits, and predicted donor codes cause the right
state-dependent answer 97.50\% of the time (random controls: 12.3\%).  Extra
rate supplies representational slack; success depends on closure over semantic
classes, not surface invariance.

\section*{\fontsize{9}{10}\selectfont Claim boundary and next tests}
\vspace{-0.3em}
The current result shows a learned, rate-limited, causally used computational
interface; it does not yet show native internal state formation or broad
reasoning.  The decisive next tests are: (i) preserve decimal-to-code entry
while training closure; (ii) hold out operation pairs and scale state entropy
to obtain a rate--reliability curve; (iii) replace addition with register
programs and Horn-rule proof states whose next action depends on the compact
frontier; and (iv) confirm three seeds and a second 7--8B family.  This connects
algorithmic length generalization \cite{velickovic2021} to a measurable
interface law: enough relevant bits, a shared decoder, and closure under every
legal update.

\vspace{-0.35em}
\begin{thebibliography}{4}
\scriptsize
\setlength{\itemsep}{0pt}
\bibitem{shannon1948} C. Shannon. A Mathematical Theory of Communication.
\emph{BSTJ}, 1948.
\bibitem{tishby2000} N. Tishby, F. Pereira, W. Bialek. The Information
Bottleneck Method. 2000.
\bibitem{shalizi2001} C. Shalizi, J. Crutchfield. Computational Mechanics:
Pattern and Prediction, Structure and Simplicity. \emph{JSP}, 2001.
\bibitem{velickovic2021} P. Veli\v{c}kovi\'c, C. Blundell. Neural Algorithmic
Reasoning. 2021.
\end{thebibliography}
\end{document}
